\documentclass[11pt]{article}
\usepackage[a4paper,margin=1in]{geometry}
\usepackage{amsmath}
\usepackage{enumitem}
\usepackage{verbatim}
\usepackage{array}

\title{\textbf{Mini Project 2 — MiniChess AI}}
\author{Student ID: 114062317}
\date{June 2026}

\begin{document}
\maketitle

\noindent\textit{This report explains the design of my MiniChess engine: the
state value function and the search algorithms (Minimax $\to$ Alpha-Beta $\to$
PVS) together with the optimizations layered on top (move ordering, quiescence,
transposition table, null-move pruning, late move reductions, aspiration
windows). It is my reference for the demo.}

\section{Overview}
The engine is a game-agnostic search core talking to a MiniChess \texttt{State}
(6$\times$5 board, king-capture win, pawns promote to Queen). My submission
policy (\texttt{submission.cpp}, a copy of my \texttt{pvs} policy) implements the
search; the state value function lives in \texttt{state.cpp::evaluate}.

The decision pipeline is the classic one:
\begin{enumerate}[noitemsep]
  \item \textbf{Iterative deepening} searches depth $1,2,3,\dots$ until the 10\,s
        budget runs out, always keeping the best move from the last completed depth.
  \item At each depth, \textbf{PVS} (an Alpha-Beta refinement) explores the tree.
  \item At the leaves, \textbf{quiescence} resolves pending captures, then the
        \textbf{evaluation} scores the position.
\end{enumerate}
All code uses only the C++ standard library --- no third-party libraries, no
inline assembly, no multithreading, no SIMD/AVX in the policy.

\section{State value function (\texttt{evaluate})}
The score is returned \textbf{from the side-to-move's perspective} (positive =
good for the player to move). A won position returns \texttt{P\_MAX}. Otherwise
\[
\text{score} = (\text{my material} + \text{my positional}) - (\text{opponent's})
              + \text{mobility} + \text{tempo}.
\]
\begin{itemize}[noitemsep]
  \item \textbf{Material} (fine $\times 10$ scale): P=20, R=60, N=70, B=80,
        Q=200, K=1000.
  \item \textbf{Piece-Square Tables} --- positional bonus per (piece, square),
        mirrored per color.
  \item \textbf{King tropism} --- bonus for pieces near the enemy king.
  \item \textbf{Mobility} --- $2\times(\text{my moves}-\text{opponent moves})$.
  \item \textbf{Passed pawns} (my addition) --- bonus growing toward promotion
        $\{0,70,45,28,16,8\}$; promotion makes a Queen, decisive on a small board.
  \item \textbf{Tempo} (my addition) --- $+6$ for the side to move.
\end{itemize}

\section{Search}

\subsection{Minimax (negamax) --- baseline}
Both sides play optimally: I maximize, the opponent minimizes. The
\textbf{negamax} identity ($\text{score}=-\,\text{opponent's score}$) lets one
routine handle both sides. This is my correctness oracle --- every faster
algorithm must return the \emph{same} score at the same depth. Winning terminals
score $\texttt{P\_MAX}-\text{ply}$ to prefer faster wins.

\subsection{Alpha-Beta pruning}
A window $[\alpha,\beta]$: $\alpha$ is what I'm already guaranteed, $\beta$ what
the opponent is. When $\alpha\ge\beta$ the opponent would never allow this line,
so the remaining moves are \textbf{pruned}. Same result as minimax, far fewer
nodes (start position, depth 6: $621{,}470 \to 26{,}392$ nodes, $\sim$23$\times$).

\subsection{Move ordering}
Best to worst: (1) transposition-table move; (2) captures by \textbf{MVV-LVA};
(3) \textbf{killer moves} (quiet moves that recently caused a cutoff at this
ply); (4) \textbf{history heuristic} (quiet moves ranked by a
$[\text{side}][\text{from}][\text{to}]$ table incremented by $\text{depth}^2$ on
cutoffs); plus a promotion bonus.

\subsection{Quiescence search}
To avoid the horizon effect, the leaf keeps searching \textbf{captures only}
until quiet. A \textbf{stand-pat} (static eval) bounds it so we never force a bad
capture. Raises selective depth from 6 to $\sim$18 at nominal depth 6.

\subsection{PVS (Principal Variation Search)}
The first (best-ordered) move is searched with the full window; later moves are
probed with a \textbf{null window} $[\alpha,\alpha{+}1]$ and only re-searched
fully if the probe beats $\alpha$. Same result as Alpha-Beta, fewer nodes.

\subsection{Transposition table}
A \textbf{Zobrist}-keyed table ($2^{22}$ entries) stores each position's score,
bound (exact/lower/upper), depth, and best move. A sufficiently deep entry cuts
off immediately; otherwise its move is tried first. Mate scores are stored
ply-relative; replacement is depth-preferred; the full 64-bit key guards against
collisions. Depth 6: $26{,}392 \to 9{,}954$ nodes; in 3\,s, depth $9\to 11$.

\subsection{Null-move pruning}
If passing (letting the opponent move twice) still beats $\beta$ after a reduced
($R=2$) search, the position is almost certainly a cutoff and is pruned. Guards:
not in check, not in pawn-only endgames (zugzwang), not on mate windows, no two
null moves in a row.

\subsection{Late Move Reductions (LMR)}
Late, quiet moves are searched at reduced depth first and only re-searched fully
if they beat $\alpha$. In a 3\,s budget this raised the reachable depth from 12
to 14.

\subsection{Aspiration windows}
Each deepening pass uses a narrow window around the previous score
($\text{prev}\pm30$), widening on a fail --- always correctness-safe. Depth
$14\to15$ in 3\,s.

\subsection{Iterative deepening \& time control}
The protocol layer drives depths $1,2,3,\dots$, emitting the best move each
depth and stopping at $\sim$half the move-time. The transposition table and the
previous PV make each pass start from the best-known move. The search is
abortable, so we never exceed 10\,s nor return an illegal move.

\section{Results}
Depth reached from the start position in a 3\,s budget:
\begin{center}
\begin{tabular}{|l|c|}
\hline
Configuration & Depth \\ \hline
Alpha-Beta + quiescence & 9 \\
+ transposition table & 11 \\
+ null-move & 12 \\
+ LMR & 14 \\
+ aspiration & 15 \\ \hline
\end{tabular}
\end{center}

Against the provided opponents (2\,s/move, 4 games, both colors):
\begin{center}
\begin{tabular}{|l|l|}
\hline
Opponent & Result (my engine) \\ \hline
\texttt{minimax-weak} & beats (4--0 at depth 6) \\
\texttt{minimax-strong} & beats (3--1 at 2\,s; 2--0 at depth 7) \\
\texttt{boss-ubgi} & \textbf{sweeps 4--0, both colors} \\ \hline
\end{tabular}
\end{center}
\textit{(Self-play tuning results to be added after the tuning phase.)}

\section{Constraints compliance}
Standard library only (no third-party code, inline asm, threads, or AVX in the
policy); $\le 10$\,s/move via abortable iterative deepening; $\le 4$\,GB
(transposition table $\sim$80\,MB); never an illegal move.

\section{Build \& run}
\begin{verbatim}
mingw32-make all            # engine + benchmark + unit tests
./build/minichess-ubgi.exe  # UBGI engine on stdin/stdout
\end{verbatim}
The submission policy is self-contained: it compiles with
\texttt{state.cpp}/\texttt{state.hpp} alone and exposes the standard
\texttt{search(State*, depth, GameHistory\&, SearchContext\&)} interface used by
the benchmark.

\end{document}
